\phantom{空一行}
\vspace{-2em} % 修正
\begin{center}
	\zihao{-2}\heiti
	基于多模型融合的红茶发酵程度识别算法研究
\end{center}

\begin{center}
	\zihao{3}\heiti
	\phantomsection   % 锚点，编译后的pdf可以直接跳转
	摘要
\end{center}

\vspace{0.5em} % center自带一定空行

这里是摘要，请好好写。本文的参考文献用于举例参考文献引用方法\cite{DP图像分类综述}。本文的附录用于举例附录引用方法如“见附录\ref{matlab_example}和附录\ref{python_example}”。

\vspace{1em}  

\noindent
\textbf{关键字：}关键字1\quad 关键字2\quad 关键字3
\newpage

\phantom{空一行}
\vspace{-2em} % 向上修正页眉与正文之间的统一空行1em
\begin{center}
	\zihao{-2}\bfseries
	\rmfamily Research on a Multi‑Model Fusion Algorithm for Black Tea Fermentation Degree Recognition
\end{center}

\begin{center}
	\phantomsection
	\zihao{3}\rmfamily\bfseries
	ABSTRACT
\end{center}

\vspace{0.5em}  % center自带一定空行

This is abstract. Literature are for example\cite{DP图像分类综述}.The appendix of this article is used for example, like this "See Appendix \ref{matlab_example} for details".

\vspace{1em}

\noindent\textbf{Key Words:}\quad Key1\quad Key2\quad Key3\quad Key4
